\clearpage
\lhead{\emph{Propuesta de experimentación}}
\chapter{Propuesta de experimentación numérica}
\label{sec:propuesta_experimental}

En este capítulo se describirá el diseño experimental empleado para comparar representaciones textuales en la predicción de resultados de licitaciones públicas paraguayas. Primero se presentarán el ambiente computacional y la procedencia de los datos. Luego se detallarán la descarga de PBC, la extracción de texto, la generación de embeddings, la definición de la etiqueta y los modelos. Finalmente se establecerán la partición, las métricas y las reglas de reproducibilidad que se utilizarán para evaluar cada alternativa.

La descripción se fundamenta en el código versionado del proyecto y en el embudo de artefactos observado. Cuando un detalle de infraestructura no quedó preservado en el repositorio, se indicará expresamente en vez de reconstruirlo mediante suposiciones.

\section{Ambiente computacional}

\subsection{Hardware}

La generación de embeddings se ejecutó con aceleración CUDA. El código selecciona una GPU cuando está disponible, utiliza \texttt{bfloat16} en ese caso y conserva los tokens en CPU hasta transferir cada microlote. La evidencia de ejecución registra 671 fallos \texttt{OutOfMemoryError}, pero el modelo exacto de GPU, su memoria, el procesador y la memoria RAM no se encuentran versionados en los artefactos disponibles. En consecuencia, no se atribuirán especificaciones de hardware no verificables.

Para reproducir el experimento será necesario registrar, además del código, el modelo de GPU, versión del controlador, versión de CUDA, memoria disponible y sistema operativo. Estos datos son relevantes porque determinan el tamaño de microlote que puede procesarse y la probabilidad de fallos en documentos extensos.

\begin{table}[H]
\centering
\caption{Información de hardware respaldada por los artefactos del proyecto.}
\label{tab:hardware}
\begin{tabular}{p{4.2cm}p{8.2cm}}
\toprule
\textbf{Componente} & \textbf{Evidencia disponible} \\
\midrule
Acelerador & GPU compatible con CUDA utilizada para generar embeddings \\
Precisión del LLM & \texttt{bfloat16} cuando CUDA está disponible \\
Gestión de memoria & Microlotes CPU--GPU, caché KV desactivada y reintentos reduciendo el microlote \\
Modelo y memoria exacta & No preservados en el repositorio; deben registrarse en una reproducción \\
Almacenamiento de artefactos & DigitalOcean Spaces mediante API compatible con S3 \\
\bottomrule
\end{tabular}
\end{table}

\subsection{Software}

El proyecto requiere Python 3.10 o superior y declara sus dependencias en \texttt{pyproject.toml}; \texttt{uv.lock} conserva resoluciones concretas por plataforma. La Tabla~\ref{tab:software} presenta componentes principales observados en el archivo de bloqueo actual. Las versiones efectivamente usadas en la corrida original deben verificarse contra el entorno archivado si se dispone de él, porque el archivo puede actualizarse después de una ejecución.

\begin{table}[H]
\centering
\caption{Componentes principales del ambiente de software versionado.}
\label{tab:software}
\begin{tabular}{p{4.4cm}p{4cm}p{4cm}}
\toprule
\textbf{Componente} & \textbf{Uso} & \textbf{Versión resuelta observada} \\
\midrule
Python & ETL y entrenamiento & $\geq 3.10$ \\
PyTorch & Embeddings y redes & 2.11.0 \\
Transformers & Carga del LLM/tokenizador & 5.4.0 \\
scikit-learn & TF--IDF, CV y métricas & 1.7.2/1.8.0 según plataforma \\
MLflow & Seguimiento experimental & 3.10.1 \\
pdfplumber & Extracción de texto PDF & 0.11.9 \\
Camelot & Extracción de tablas & 1.0.9 \\
pandas & Manipulación tabular & 2.3.3 \\
\bottomrule
\end{tabular}
\end{table}

\section{Fuente y procedencia de los datos}

\subsection{API de la DNCP}

Los procesos provienen del servicio de datos abiertos de la DNCP. El script \texttt{fetch\_ids\_by\_status.py} consulta el endpoint:

\begin{center}
\small\url{https://www.contrataciones.gov.py/datos/api/v3/doc/search/processes}
\end{center}

La consulta utiliza \texttt{items\_per\_page=100}, el código de clasificación \texttt{72131701}, el estado de la licitación y el número de página. El endpoint forma parte de la publicación paraguaya basada en Open Contracting \cite{dncpOpenContracting,dncpDatosAbiertos}. Se recorrieron las páginas disponibles para \texttt{unsuccessful} y \texttt{cancelled}; cada respuesta se normalizó añadiendo el estado utilizado en la consulta.

De cada registro se recuperaron, cuando estaban disponibles, el OCID, el \texttt{tender.id} y los identificadores de adjudicaciones. Si no aparecía \texttt{tender.id}, el código utilizó el OCID como alternativa. El identificador de tender se mantuvo como clave para unir etapas posteriores.

\subsection{Conjunto inicial}

El conjunto inicial reunió tres fuentes almacenadas en DigitalOcean Spaces:

\begin{enumerate}
  \item \texttt{procurements.json}, con 10.000 procesos materializados como \texttt{complete} desde una instantánea previa de la DNCP;
  \item \texttt{ids\_unsuccessful.json}, con 257 procesos recuperados de las tres páginas disponibles;
  \item \texttt{ids\_cancelled.json}, con 371 procesos recuperados de las cuatro páginas disponibles.
\end{enumerate}

La unión produjo 10.628 filas. El orden de combinación fue no exitosos, cancelados y completos. Ante identificadores duplicados, la primera capa conservó prioridad y las siguientes rellenaron campos ausentes. El mismo orden definió la cola de descarga, con el propósito de procesar primero la clase minoritaria.

\section{Conjunto de datos}
La instantánea de metadatos utilizada en este trabajo se construyó a partir de registros de contratación pública publicados por la Dirección Nacional de Contrataciones Públicas (DNCP) de Paraguay y recuperados mediante el endpoint oficial de datos abiertos \texttt{v3/doc/search/processes}. El conjunto inicial contiene 10.628 contrataciones de construcción vial bajo el código de catálogo \texttt{72131701}: 10.000 procesos completados materializados por una canalización previa de extracción de la DNCP y 628 procesos adicionales no exitosos recolectados según su estado. Para reducir el desbalance de clases, la recuperación y el procesamiento priorizaron las contrataciones no exitosas antes que las completadas.

El subconjunto supervisado es menor porque se requieren varios pasos de filtrado. La mayor reducción ocurre antes del procesamiento textual: solo se recuperaron 1.845 PDF de PBC, principalmente porque los límites de solicitudes de la DNCP impidieron realizar una búsqueda y descarga exhaustivas dentro del tiempo disponible. Luego, la canalización extrajo texto de 1.842 PDF y generó embeddings congelados del LLM para 1.162 documentos. Los tres fallos de extracción correspondieron a documentos escaneados no compatibles con el extractor sin OCR, mientras que la mayoría de los fallos de generación de embeddings fueron errores por falta de memoria en la GPU, después de una única ejecución limitada por el costo de infraestructura. Por ello, las conclusiones se aplican a este subconjunto procesado y no automáticamente a toda la población de contrataciones.

La distribución de clases del subconjunto procesado presenta un desbalance moderado: 844 contrataciones ($72.63\%$) están etiquetadas como exitosas (\texttt{complete}) y 318 ($27.37\%$) como no exitosas (\texttt{unsuccessful}/\texttt{cancelled}/\texttt{canceled}). La longitud media representada es de 7,46 fragmentos por contratación, con una desviación estándar de 13,52, una mediana de 3 fragmentos y un máximo de 245. Estos valores indican una distribución de longitudes muy asimétrica, con numerosos documentos cortos y una cantidad menor de documentos considerablemente más largos. La dimensión de los embeddings heredada del LLM base congelado es 2880.

El subconjunto procesado también está menos desbalanceado que la población general de contrataciones. En la instantánea completa, los resultados exitosos son mucho más frecuentes: aproximadamente 10.000 contrataciones exitosas frente a unas 650 no exitosas. Por consiguiente, las métricas dependientes del umbral y las interpretaciones sensibles a la prevalencia deben considerarse válidas para el subconjunto procesado, no como puntos operativos directamente desplegables.

\section{Tratamiento de la fuga de información}
El modelo opera exclusivamente con información disponible en el momento de la inferencia. Todas las entradas se derivan de documentos de licitación producidos y publicados durante la etapa de convocatoria, antes de que se conozca el resultado final de la contratación. Esto hace que la configuración predictiva sea operativamente realista y reduce sustancialmente el riesgo de fuga directa a partir de contenido explícito sobre el estado final.

En consecuencia, la principal amenaza restante para la validez no es una fuga evidente y posterior de la variable objetivo, sino las correlaciones predictivas indirectas y la representatividad limitada. Los documentos de licitación pueden codificar regularidades estructurales, institucionales o estilísticas que predicen el resultado sin ser causales. Estas regularidades son legítimas para la predicción ex ante porque están disponibles en el momento de la decisión, pero no deben interpretarse como causas directas del éxito o fracaso de una contratación.

\section{Protocolo de evaluación}
El modelo se evalúa mediante validación cruzada estratificada de cinco pliegues. El subconjunto procesado de 1.162 ejemplos se divide en cinco pliegues que preservan las etiquetas, de modo que cada ronda utiliza aproximadamente cuatro quintas partes de los datos para ajustar el modelo y una quinta parte como partición de evaluación fuera de pliegue. En términos concretos, cada ronda entrena con unas 930 contrataciones y evalúa aproximadamente 232 que no se utilizaron para ajustar el modelo de esa ronda.

La estratificación es importante porque la clase positiva no está balanceada: cerca del $72.6\%$ del subconjunto procesado corresponde a resultados exitosos. Conservar esta prevalencia entre pliegues reduce variaciones evitables en F1, exactitud balanceada y PR AUC. Los valores finales informados son la media y la desviación estándar de los cinco pliegues reservados. Para el análisis de sensibilidad al umbral, los cuatro pliegues de entrenamiento de cada ronda se dividen nuevamente en una porción interna de entrenamiento y otra de validación; el umbral que maximiza F1 en la validación interna se aplica después al pliegue de evaluación externo de esa ronda.

\section{Métricas}
La evaluación utiliza las siguientes métricas:

\begin{itemize}
  \item \textbf{ROC AUC}, para medir la calidad del ordenamiento con independencia del umbral de decisión.
  \item \textbf{PR AUC} (precisión promedio), para evaluar el ordenamiento con atención explícita a la clase positiva.
  \item \textbf{Puntaje F1}, con un umbral de probabilidad de $0.5$.
  \item \textbf{Exactitud balanceada}, también con un umbral de $0.5$.
  \item \textbf{Puntaje de Brier}, para medir la calidad de la calibración probabilística.
  \item \textbf{Pérdida logarítmica}, como regla de puntuación propia sobre las probabilidades predichas.
\end{itemize}

La elección de métricas ante el desbalance depende del objetivo operativo: los resúmenes basados en PR enfatizan la recuperación de la clase positiva, mientras que ROC AUC y PR AUC responden de manera distinta a la prevalencia y deben interpretarse en contexto \cite{richardson2024roc}. Debido a que PR AUC puede inducir a error cuando la clase positiva es frecuente, se interpreta junto con la prevalencia positiva $\pi$ en los pliegues reservados. La implementación también registra \texttt{positive\_prevalence} y \texttt{pr\_auc\_excess\_over\_prevalence}; esta última se define como PR AUC menos $\pi$. La cantidad ofrece una estimación compacta de cuánto mejora el modelo sobre la frecuencia de la clase positiva, utilizada habitualmente como referencia al discutir el comportamiento PR sin capacidad predictiva en contextos desbalanceados.


\section{Recuperación de documentos PBC}

Para cada proceso, \texttt{merge\_and\_download\_pbcs.py} consultará la información documental asociada e intentará localizar un PBC o carta de invitación. La ejecución documentada aceptó PDF como artefacto final. El código inspecciona título, URL y formato; también contiene rutas auxiliares para archivos comprimidos o documentos de oficina, pero el conjunto final observado está compuesto por PDF.

Los archivos se almacenan bajo la clave lógica:

\begin{center}
\texttt{outcome-predictor/pbcs/pdf/\{tenderId\}.pdf}
\end{center}

El proceso es reanudable: antes de descargar, puede comprobar si el objeto existe. El conjunto JSON enriquecido conserva indicadores de descarga, clave remota y, cuando se registró, motivo de omisión o error. Una reconciliación final compara las filas con el listado real de objetos en Spaces para evitar que el JSON declare un PDF inexistente.

Se obtuvieron 1.845 PDF. Esta cifra representa los archivos recuperados por la canalización, no todos los PBC existentes en la DNCP. Para 42 procesos se conservó un motivo de fallo; 8.741 filas no contienen evidencia de intento o causa final y se documentan como no procesadas dentro del tiempo disponible, principalmente por limitación de solicitudes. Esta diferencia exige evitar tasas de disponibilidad calculadas como si toda la población hubiese sido consultada exhaustivamente.

\section{Extracción y normalización del texto}

Los PDF descargados se procesaron con \texttt{PDFReader}. La etapa utiliza \texttt{pdfplumber} para texto y puede utilizar Camelot en modos \textit{lattice} y \textit{stream} para tablas. Cada resultado se almacenó en:

\begin{center}
\texttt{outcome-predictor/pbcs/txt/\{tenderId\}.txt}
\end{center}

De 1.845 PDF se obtuvieron 1.842 textos. Los tres fallos correspondieron a documentos escaneados para los que el extractor no incorporaba OCR. No se imputó texto vacío ni se reemplazaron esos documentos por contenido de otra fuente; quedaron excluidos.

La extracción puede conservar ruido de paginación o alterar tablas. El experimento tratará el texto obtenido como entrada efectiva del modelo y no como transcripción perfecta del PBC. Una extensión con OCR deberá evaluarse por separado porque cambiaría la población incluida y la calidad de la entrada.

\section{Generación de embeddings}

\subsection{Codificador congelado}

El script \texttt{embed\_pbcs.py} cargará \texttt{openai/gpt-oss-20b} mediante Hugging Face Transformers y congelará todos sus parámetros. La configuración seleccionará \texttt{bfloat16} en CUDA, desactivará la caché de claves y valores y utilizará el backbone que devuelve el último estado oculto.

La secuencia completa se tokenizará sin truncamiento. Se dividirá en fragmentos de 4096 tokens con desplazamiento de 2048. Cada fragmento se completará hasta la longitud máxima para formar microlotes. El embedding corresponderá al estado oculto del último token válido, no al relleno.

\subsection{Control de memoria}

El tamaño inicial del microlote será cuatro. Ante un error de memoria CUDA, el código vaciará la caché, reducirá el tamaño por mitades y reintentará hasta llegar a uno. Si el error persiste, se registrará y el documento quedará sin embedding. Este procedimiento evita que un fallo detenga toda la colección, aunque induce una selección técnica por longitud y demanda de memoria.

Cada archivo \texttt{.pt} conserva \texttt{tender\_id}, matriz de embeddings, identificador del modelo, longitud, desplazamiento, tamaño de microlote y eventual etiqueta. Se produjeron 1.162 archivos. Entre los 680 textos sin embedding se observaron 671 errores de memoria y nueve \texttt{ValueError}.

\section{Embudo de construcción del conjunto supervisado}

La Tabla~\ref{tab:funnel} presenta el recorrido completo. Los conteos provienen de objetos observados y del JSON consolidado.

\begin{table}[H]
\centering
\caption{Embudo de datos desde procesos DNCP hasta el conjunto supervisado.}
\label{tab:funnel}
\begin{tabular}{p{5.3cm}rrr}
\toprule
\textbf{Etapa} & \textbf{Universo} & \textbf{Resultado válido} & \textbf{Pérdida} \\
\midrule
Procesos unificados & 10.628 & 10.628 & 0 \\
PBC PDF recuperados & 10.628 & 1.845 & 8.783 \\
Texto extraído & 1.845 & 1.842 & 3 \\
Embeddings generados & 1.842 & 1.162 & 680 \\
Ejemplos supervisados & 1.162 & 1.162 & 0 \\
\bottomrule
\end{tabular}
\end{table}

La Tabla~\ref{tab:funnel_status} desagrega etapas por estado. La priorización elevó la proporción de no exitosos respecto de la instantánea completa.

\begin{table}[H]
\centering
\caption{Conteos por estado en etapas principales.}
\label{tab:funnel_status}
\begin{tabular}{lrrrr}
\toprule
\textbf{Estado} & \textbf{Inicial} & \textbf{PDF} & \textbf{Texto} & \textbf{Embedding} \\
\midrule
\texttt{complete} & 10.000 & 1.273 & 1.270 & 844 \\
\texttt{unsuccessful} & 257 & 246 & 246 & 204 \\
\texttt{cancelled} & 371 & 326 & 326 & 114 \\
\midrule
Total & 10.628 & 1.845 & 1.842 & 1.162 \\
\bottomrule
\end{tabular}
\end{table}

La pérdida no es aleatoria. Se priorizaron estados minoritarios y los errores de embedding afectaron principalmente documentos costosos. Por ello, el subconjunto final no debe presentarse como muestra representativa del universo paraguayo. El protocolo medirá desempeño interno sobre los 1.162 casos procesados.

\section{Definición de la instancia y etiqueta}

El estudio busca anticipar el estado con el que cada contratación queda registrada a partir del texto de su PBC. Para ello se utilizan únicamente procesos que cuentan con un documento aprovechable y con uno de los estados finales considerados. Los estados se agrupan en dos resultados:

\begin{itemize}
  \item Se considera exitoso el proceso registrado como \texttt{complete}.
  \item Se considera no exitoso el proceso registrado como \texttt{unsuccessful}, \texttt{cancelled} o \texttt{canceled}.
\end{itemize}

Los términos \texttt{cancelled} y \texttt{canceled} expresan el mismo estado, ``cancelado'', con las grafías británica y estadounidense del inglés; ambos aparecen en los datos de origen y por eso se normalizan dentro del mismo grupo. Reducir los estados a estas dos alternativas convierte el problema en una clasificación binaria: para cada PBC, el modelo estima la probabilidad de que el proceso pertenezca al primer grupo.

Los registros con otros estados se excluyen del entrenamiento y de la evaluación. Esta agrupación es una decisión operativa basada en la fuente disponible. Un proceso \texttt{complete} no implica necesariamente máximo valor público, y una cancelación puede responder a causas legítimas; por tanto, la etiqueta no constituye una evaluación integral de éxito, eficiencia o integridad.


El mapeo se implementa en \texttt{src/data/procurement\_target.py} y se alinea mediante \texttt{tenderId}. Registros sin uno de los estados admitidos se excluirán. La etiqueta no se extraerá del PBC y no se suministrará al codificador.

\section{Modelos que se evaluarán}

\section{Representación documental}
Cada documento PBC se convierte primero de PDF a texto mediante la canalización de extracción previa. El texto resultante se tokeniza y divide en fragmentos superpuestos. En la implementación actual, cada fragmento tiene una longitud de 4096 tokens y un desplazamiento de 2048 tokens. Se utiliza el modelo de lenguaje causal preentrenado \texttt{openai/gpt-oss-20b} como codificador documental congelado \cite{agarwal2025gptoss}. Este modelo se seleccionó como un LLM práctico de pesos abiertos para extraer representaciones densas de textos de contratación en español y, a la vez, mantener reproducible la canalización experimental. Congelar el codificador también evita ajustar un modelo grande sobre un conjunto de datos modesto y permite que la comparación se concentre en cuánta señal existe en las representaciones preentrenadas y cómo se agrega posteriormente, en consonancia con trabajos recientes sobre embeddings textuales basados en LLM \cite{wang2024improvingembeddings}. Los fragmentos se procesan en microlotes para reducir la presión sobre la memoria de la GPU.

Para cada fragmento, el sistema extrae la representación oculta asociada al último token válido. La salida correspondiente a un documento de contratación es, por tanto, una secuencia de embeddings de fragmentos:
\[
\mathbf{E} = [\mathbf{e}_1, \mathbf{e}_2, \dots, \mathbf{e}_N], \qquad \mathbf{e}_i \in \mathbb{R}^{d_{in}},
\]
donde $N$ es la cantidad de fragmentos del documento y $d_{in}$ está determinado por el tamaño oculto del modelo preentrenado. Esta configuración permite comparar dos estrategias de representación basadas en LLM: un vector documental simple obtenido mediante el promedio de los fragmentos y un modelo secuencial que aprende las interacciones entre ellos.

\section{Arquitectura del predictor}
El clasificador principal, \texttt{TenderSuccessPredictor}, es un transformer liviano aplicado sobre los embeddings de los fragmentos. Sea $\mathbf{E} \in \mathbb{R}^{N \times d_{in}}$ la secuencia de fragmentos de un documento. El modelo aplica:

\begin{enumerate}
  \item Una proyección lineal de $d_{in}$ a una dimensión latente entrenable $d_{model}$.
  \item La inserción de un token aprendido \texttt{[CLS]} al comienzo de la secuencia.
  \item Un codificador transformer sobre la secuencia completa de vectores de fragmentos.
  \item Una pequeña cabeza de clasificación multicapa sobre la representación final de \texttt{[CLS]}.
\end{enumerate}

En los experimentos actuales, los hiperparámetros posteriores son: $d_{model}=128$, $n_{heads}=4$, dimensión de la red prealimentada $=512$, dropout $=0.3$ y una capa de codificación. Se emplean máscaras de relleno para agrupar de forma segura documentos con diferentes cantidades de fragmentos.

El modelo produce un único logit por documento. Una transformación sigmoide convierte este logit en una probabilidad de éxito:
\[
\hat{p}(y=1 \mid \mathbf{E}) = \sigma(z).
\]
El entrenamiento utiliza entropía cruzada binaria con logits.

\section{Líneas base}
Para interpretar los resultados basados en LLM, se evalúan tres clases de líneas base.

En primer lugar, se emplean dos referencias triviales: un clasificador de clase mayoritaria y un clasificador aleatorio cuya tasa positiva coincide con la prevalencia empírica. Estas líneas base establecen la región esperada sin capacidad predictiva, donde ROC AUC debería situarse cerca de $0.5$ y PR AUC cerca de la prevalencia positiva.

En segundo lugar, se utiliza una línea base léxica con características TF--IDF de hasta 50.000 unigramas y bigramas, frecuencia documental mínima de $5$ y un clasificador de regresión logística con pesos de clase balanceados. Se utiliza TF--IDF en lugar de embeddings densos estáticos como Word2Vec porque proporciona una referencia sólida e interpretable de bolsa de palabras que sigue siendo competitiva en documentos extensos \cite{wang2012baselines,alvaprincipe2025survey}. Además, permite aislar la comparación con representaciones congeladas de LLM sin añadir otra etapa entrenable de embeddings bajo supervisión limitada. Las líneas base lineales de bolsa de palabras son referencias habituales en clasificación de textos \cite{wang2012baselines}, y la evidencia reciente sobre documentos extensos indica que las representaciones léxicas aún son competitivas en algunos escenarios \cite{alvaprincipe2025survey}.

En tercer lugar, se definen dos líneas base intermedias basadas en LLM que utilizan los mismos embeddings preentrenados de fragmentos que el sistema principal, pero sustituyen el transformer entre fragmentos por una agregación documental más simple. En la primera variante, se promedian los embeddings de los fragmentos y se suministran a una regresión logística. En la segunda, la misma representación promedio alimenta un pequeño perceptrón multicapa con una capa oculta. Estas líneas base aíslan el valor de la representación del valor aportado por el modelado explícito de las interacciones entre fragmentos.

\section{Entrenamiento y reproducibilidad}
El entrenamiento utiliza AdamW, una tasa de aprendizaje de $10^{-4}$, decaimiento de pesos de $0.01$, lotes de tamaño $8$ y hasta $50$ épocas. Dentro de cada ronda de validación cruzada, una época corresponde a un recorrido completo por la porción de entrenamiento de esa ronda. Después de cada época se calcula la pérdida de validación sobre la porción reservada; si no mejora durante cinco épocas consecutivas, el entrenamiento se detiene y se restaura el mejor punto de control. La reproducibilidad se controla mediante semillas fijas, comportamiento determinista de cuDNN, cargadores de datos con semilla y reinicios independientes del estado aleatorio al comienzo de cada pliegue.


\subsection{Líneas base triviales}

Se incluirá un clasificador de clase mayoritaria y uno aleatorio cuya tasa positiva coincida con la prevalencia. Servirán como control del pipeline: ROC AUC deberá situarse cerca de 0,5 y PR AUC cerca de la prevalencia, dentro de la variación esperable.

\subsection{TF--IDF con regresión logística}

El vectorizador se ajustará exclusivamente con textos de entrenamiento de cada pliegue. Se usarán 50.000 características como máximo, unigramas y bigramas, y frecuencia documental mínima de cinco. La regresión logística utilizará \texttt{class\_weight=balanced}. Transformar el pliegue externo con un vocabulario aprendido solo en entrenamiento evitará fuga por frecuencia documental.

\subsection{Promedio con regresión logística}

Se calculará el promedio de embeddings por documento y se ajustará una regresión logística. Este modelo aislará el contenido de la representación congelada sin introducir una red secuencial.

\subsection{Promedio con MLP}

El mismo promedio alimentará una red con una capa oculta. En el análisis de sensibilidad se utilizó dimensión oculta 256, ReLU, dropout 0,2, AdamW, lotes de 64 y detención temprana. La comparación con regresión logística permitirá observar si una transformación no lineal del promedio aporta valor.

\subsection{Transformer entre fragmentos}

La matriz variable de embeddings se rellenará dentro de cada lote y se acompañará con una máscara. La arquitectura proyectará de 2880 a 128 dimensiones, añadirá \texttt{[CLS]}, utilizará cuatro cabezas, dimensión prealimentada 512, dropout 0,3 y una capa. La cabeza final producirá un logit.

\section{Partición y protocolo de evaluación}

Se utilizará validación cruzada estratificada de cinco pliegues con semilla 42. Cada pliegue conservará aproximadamente la prevalencia. Los modelos comparados emplearán los mismos índices externos. El vectorizador y los clasificadores se ajustarán nuevamente en cada ronda.

Para las métricas a umbral fijo se utilizará $\tau=0.5$. Para el experimento de sensibilidad, los cuatro pliegues externos de entrenamiento se dividirán en entrenamiento interno y validación interna, con una proporción de validación de 0,2. Se seleccionará el umbral que maximice F1 en la partición interna; ante empates se preferirá el más cercano a 0,5 y luego el menor. El umbral se aplicará una sola vez al pliegue externo.

Este anidamiento separará selección y evaluación. No se utilizará el pliegue externo para elegir hiperparámetros, vocabulario o umbral.

\section{Entrenamiento}

El transformer se entrenará con AdamW, tasa $10^{-4}$, decaimiento 0,01, lotes de ocho y hasta 50 épocas. Después de cada época se calculará la pérdida de validación. Si no mejora durante cinco épocas consecutivas, se detendrá el entrenamiento y se restaurará el mejor estado.

La pérdida será entropía cruzada binaria con logits. El uso directo de logits evita aplicar sigmoide antes de una función que implementa de forma estable ambas operaciones. Durante evaluación, la sigmoide transformará logits en probabilidades.

\section{Métricas de evaluación}

Se informarán ROC AUC, PR AUC, F1, exactitud balanceada, Brier y pérdida logarítmica. ROC AUC y PR AUC evaluarán ranking; F1 y exactitud balanceada dependerán del umbral; Brier y pérdida logarítmica evaluarán probabilidades. También se registrará la prevalencia positiva y $PR\ AUC-\pi$.

Las curvas de precisión-exhaustividad se construirán agregando predicciones externas. Las curvas de calibración agruparán probabilidades y compararán predicción media con frecuencia observada. Dado el tamaño limitado, los puntos de bins extremos se interpretarán con cautela.

\section{Trazabilidad y reproducibilidad}

La semilla global será 42 y cada pliegue recibirá una semilla derivada independiente. El código configurará cuDNN determinista, generadores de \texttt{DataLoader} y reinicios de estados aleatorios. MLflow conservará hiperparámetros y métricas. Los objetos pesados residirán en Spaces y los metadatos de seguimiento, por defecto, en SQLite local.

Los artefactos públicos archivados en Zenodo incluyen código, configuración y documentación, pero excluyen datos locales, credenciales y salidas pesadas. La reproducción completa requerirá acceso legítimo a los datos fuente o una nueva descarga desde la DNCP, además del modelo base. Esta limitación deberá declararse al presentar el grado de reproducibilidad.

\section{Amenazas a la validez previstas}

Se considerarán cuatro grupos:

\begin{enumerate}
  \item \textbf{Validez de constructo}: el estado final es una aproximación operacional al éxito.
  \item \textbf{Validez interna}: podrían existir duplicados documentales, regularidades institucionales o señales indirectas; el protocolo evita datos posteriores, pero no convierte correlaciones en causalidad.
  \item \textbf{Validez externa}: el subconjunto está filtrado por disponibilidad y cómputo, y se limita a construcción vial paraguaya.
  \item \textbf{Validez de conclusión}: cinco pliegues sobre 1.162 ejemplos producen incertidumbre; diferencias pequeñas no deben presentarse como dominancia universal sin pruebas adicionales.
\end{enumerate}

Estas amenazas guiarán la interpretación del capítulo de resultados y evitarán extrapolar los umbrales del subconjunto a la distribución completa.
